\documentclass[a5paper,10pt]{article}

% https://github.com/InDevRus/filippov-solutions

% Нумерация страниц
\pagenumbering{gobble}

% Настройка полей
\usepackage{geometry}
\geometry{left=1cm}
\geometry{right=1cm}
\geometry{top=1cm}
\geometry{bottom=1.5cm}

% Вставка таблицы
\usepackage{array}
\usepackage{tabularx}

% Инструменты выравнивания формул
\usepackage{amsmath}
\usepackage{mathtools}

% Диакритика в математике
\usepackage{amsfonts}

% Вычисления размеров на ходу
\usepackage{calc}

% Удобные записи производных
\usepackage[italic=true]{derivative}

% Настройки сносок
\usepackage{enotez}

% Длинные знаки векторов
\usepackage{anyfontsize}
\usepackage[b]{esvect}

% Вставка картинок
\usepackage{graphicx}

% Вставка красной строки
\usepackage{indentfirst}
\setlength{\parindent}{1em}

% Условия в командах
\usepackage{xifthen}

% Луа-скрипты
\usepackage{luacode}

% Настройка команд отдельно в математическом и текстовом режимах
\usepackage{mathcommand}

% Для повторения LaTeX кода
\usepackage{multido}

% Конфигурация отступов формул
\delimitershortfall-1sp
\usepackage{mleftright}
\mleftright

% Растягивание объектов
\usepackage{relsize}

% Обтекание текстом
\usepackage{wrapfig}
\usepackage{subcaption}
\usepackage{floatflt}

% Поддержка ввода кириллицы
\usepackage[english, russian]{babel}

% Настройка корневой позиции для компилятора
\usepackage{import}
\providecommand*{\ProjectRootPath}{..}

\import{\ProjectRootPath/preambles/macros/}{theorems}

\import{\ProjectRootPath/preambles/fonts}{font_setup}

% Знак градуса
\usepackage{siunitx}

\import{\ProjectRootPath/preambles/macros/}{operators}
\import{\ProjectRootPath/preambles/macros/}{redefinitions}

\import{\ProjectRootPath/preambles/fonts}{letters_setup}
\import{\ProjectRootPath/preambles/fonts/adjustments}{custom_superscripts}

\import{\ProjectRootPath/preambles/custom}{formatting}
\import{\ProjectRootPath/preambles/custom}{frames}
\import{\ProjectRootPath/preambles/custom}{list_setup}

% TODO: перевести команды на современный стиль объявления

\begin{document}

$\tsys{
    P\br{x; y} = \pow{x}{2} - y
    \\ Q\br{x; y} = \pow{x}{2} - \br{y - 2}^2
}$;
$\tsys{
    y = \pow{x}{2}
    \\ \pow{x}{2} - \br{\pow{x}{2} - 2}^2 = 0
}$;
$\br{x - \br{\pow{x}{2} - 2}}\br{x + \br{\pow{x}{2} - 2}} = 0$. \\
Особыми точками будут 
$\sub{A}{1} = \br{-2; 4}$,
$\sub{A}{2} = \br{-1; 1}$,
$\sub{A}{3} = \br{1; 1}$,
$\sub{A}{4} = \br{2; 4}$.

\begin{enumerate}
    \item $\tsys{x = \xi - 2 \\ y = \phi + 4}$;
    $\tsys{
        \dot{\xi} = -4\xi - \phi + \pow{\xi}{2} \\
        \dot{\phi} = -4\xi - 4\phi + \pow{\xi}{2} - \pow{\phi}{2}
    }$;
    $A = \begin{pmatrix} -4 & -1 \\ -4 & -4 \end{pmatrix}$.
    $\br{\lambda + 4}^2 - 4 = 0$.
    $\sub{\lambda}{1} = -6$, $\sub{\lambda}{2} = -2$.
    Это значит, что особая точка $\sub{A}{1}\br{-2; 4}$ -- <<узел>>, причём устойчивый.
    
    $A - \sub{\lambda}{1} E \sim \begin{pmatrix} 2 & -1 \end{pmatrix}$,
    $\sub{v}{1} = \begin{pmatrix} 1 \\ 2 \end{pmatrix}$.
    
    $A - \sub{\lambda}{2} E \sim \begin{pmatrix} 2 & 1 \end{pmatrix}$,
    $\sub{v}{2} = \begin{pmatrix} 1 \\ -2 \end{pmatrix}$.
    
    $\sub{\phi}{1} = 2\xi$ -- трансверсаль, $\sub{\phi}{2} = -2\xi$ -- касательная.
    
    \item $\tsys{x = \xi - 1 \\ y = \phi + 1}$;
    $\tsys{
        \dot{\xi} = -2\xi - \phi + \pow{\xi}{2} \\
        \dot{\phi} = -2\xi + 2\phi + \pow{\xi}{2} - \pow{\phi}{2}
    }$;
    $A = \begin{pmatrix} -2 & -1 \\ -2 & 2 \end{pmatrix}$.
    $\pow{\lambda}{2} - 4 - 2 = 0$.
    $\sub{\lambda}{1} = -\sqrt{6}$,
    $\sub{\lambda}{2} = \sqrt{6}$,
    что означает, что особая точка $\sub{A}{2} = \br{-1; 1}$ -- <<седло>>.
    
    $A - \sub{\lambda}{1} E \sim \begin{pmatrix} 2 - \sqrt{6} & 1 \end{pmatrix}$,
    $\sub{v}{1} = \begin{pmatrix} 1 \\ -2 + \sqrt{6} \end{pmatrix}$.
    
    $A - \sub{\lambda}{2} E \sim \begin{pmatrix} 2 + \sqrt{6} & 1 \end{pmatrix}$,
    $\sub{v}{2} = \begin{pmatrix} 1 \\ -2 - \sqrt{6} \end{pmatrix}$.
    
    $\sub{\phi}{1} = \br{-2 + \sqrt{6}} \xi$ 
    и $\sub{\phi}{2} = \br{-2 - \sqrt{6}} \xi$ -- сепаратрисы.
    
    \item $\tsys{x = \xi + 1 \\ y = \phi + 1}$;
    $\tsys{
        \dot{\xi} = 2\xi - \phi + \pow{\xi}{2} \\
        \dot{\phi} = 2\xi + 2\phi + \pow{\xi}{2} - \pow{\phi}{2}
    }$;
    $A = \begin{pmatrix} 2 & -1 \\ 2 & 2 \end{pmatrix}$.
    $\br{\lambda - 2}^2 + 2 = 0$.
    $\sub{\lambda}{1} = 2 - \sqrt{2}i$,
    $\sub{\lambda}{2} = 2 + \sqrt{2}i$.
    Таким образом, особая точка $\sub{A}{3} = \br{1; 1}$ -- <<фокус>>, причём неустойчивый.
    В точке $\tsys{\xi = 1 \\ \phi = 0}$ получаем вектор $\tsys{\dot{\xi}\br{1; 0} = 3 \\ \dot{\phi}\br{1; 0} = 3}$, так что направление раскручивания траекторий -- против часовой.
    
    \item $\tsys{x = \xi + 2 \\ y = \phi + 4}$;
    $\tsys{
        \dot{\xi} = 4\xi - \phi + \pow{\xi}{2} \\
        \dot{\phi} = 4\xi - 4\phi + \pow{\xi}{2} - \pow{\phi}{2}
    }$;
    $A = \begin{pmatrix} 4 & -1 \\ 4 & -4 \end{pmatrix}$.
    $\pow{\lambda}{2} - 16 + 4 = 0$.
    $\sub{\lambda}{1} = -2\sqrt{3}$, $\sub{\lambda}{2} = 2\sqrt{3}$,
    а особая точка $\sub{A}{4}\br{2;  4}$ -- <<седло>>.
    
    $A - \sub{\lambda}{1} E \sim \begin{pmatrix} 4 + 2\sqrt{3} & -1 \end{pmatrix}$,
    $\sub{v}{1} = \begin{pmatrix} 1 \\ 4 + 2\sqrt{3} \end{pmatrix}$.
    
    $A - \sub{\lambda}{2} E \sim \begin{pmatrix} 4 - 2\sqrt{3} & -1 \end{pmatrix}$,
    $\sub{v}{2} = \begin{pmatrix} 1 \\ 4 - 2\sqrt{3} \end{pmatrix}$.
    
    Сепаратрисами будут $\sub{\phi}{1} = \br{4 + 2\sqrt{3}} \xi$ 
    и $\sub{\phi}{2} = \br{4 - 2\sqrt{3}} \xi$.
\end{enumerate}

\end{document}
